\begin{abstract}
A map lens that \emph{summarises} the data under it---rather than magnifying or
filtering it---must decide what to summarise and where to put the summary.
Existing lens models leave both inside a single lens function and join. We
refine them into seven stages (selection, binning, normalisation, placement,
anchor, marks, association) and show that a category ring, a bearing-faithful
necklace, a route profile and a gridded glyphmap are points in one space,
reached by changing arguments. The part of the space we find new is where the
lens boundary itself carries data: aggregates placed on a ring at the bearing of
what they summarise, displaced minimally rather than dropped; a leader that
appears exactly where placement or unrolling has given adjacency away; a ring
that unrolls into an axis without re-solving; within-unit structure read from
the lens's own polar frame; and a radius control that shows how fragile the
reading is. Every demonstration summarises one categorical variable by count:
four classes of OpenStreetMap places in Yogyakarta, as 1{,}449 points and as
counts for 14 districts. We relate the model to lens design spaces, focus-region labelling,
necklace maps and geographically weighted statistics, state which parts we
concede to them, describe an open-source implementation, and set out the
controlled study that the central design claim still requires.
\todo{Rewrite once the study has run; keep to the venue's word limit.}
\end{abstract}
